\chapter{연립일차방정식과 Gauss 소거법}

\chapterintro{이 장에서는 연립일차방정식을 행렬 언어로 옮긴 뒤, 왜 소거 절차가 해를 바꾸지 않는지 엄밀하게 확인하겠습니다. 계산 알고리즘을 익히는 것에 그치지 않고, 해의 존재성, 유일성, 자유도의 구조까지 한 흐름으로 연결해 보겠습니다.}
\chaptergoals{연립일차방정식을 계수행렬과 증강행렬로 표준화할 수 있다.}{기본 행연산이 해집합을 보존함을 완전하게 설명할 수 있다.}{RREF와 랭크를 이용해 해의 구조를 매개변수 형태로 기술할 수 있다.}

\sectiontemplate{선형시스템의 행렬 표현}{연립일차방정식(linear system)은 선형대수 전체를 여는 출발점입니다. 먼저 표기를 통일해 두면 이후의 계산과 증명이 모두 짧아지고, 서로 다른 문제를 같은 언어로 비교할 수 있습니다.}{\item 선형시스템의 표준형을 정확히 쓸 수 있습니다.\item 계수행렬과 증강행렬을 구분할 수 있습니다.\item 동차/비동차 시스템의 해집합 차이를 설명할 수 있습니다.}{\item 체와 벡터 표기\item 행렬의 기본 인덱스 표기}

\begin{definition}
체 $K$와 양의 정수 $m,n$에 대하여
\[
\sum_{j=1}^n a_{ij}x_j=b_i \quad (i=1,\dots,m)
\]
꼴의 식들을 모아 놓은 것을 $K$ 위의 \emph{연립일차방정식}이라 한다. 여기서 $a_{ij},b_i\in K$이다.
\end{definition}

\begin{definition}
위 시스템에서 $A=(a_{ij})\in M_{m,n}(K)$를 \emph{계수행렬}(coefficient matrix), $\mathbf{b}=(b_1,\dots,b_m)^T\in K^m$를 상수벡터라 한다. 시스템은
\[
A\mathbf{x}=\mathbf{b}
\]
로 쓸 수 있으며, $(A\mid \mathbf{b})$를 \emph{증강행렬}(augmented matrix)이라 한다. 특히 $\mathbf{b}=\mathbf{0}$이면 \emph{동차시스템}(homogeneous system), $\mathbf{b}\neq \mathbf{0}$이면 비동차시스템이라 한다.
\end{definition}

\paragraph{증명 전략.}
동차시스템의 해집합이 부분공간임을 보이기 위해 영벡터 포함, 덧셈 닫힘, 스칼라곱 닫힘을 차례로 확인한다. 핵심 도구는 행렬곱의 선형성이다.

\begin{theorem}[동차해집합의 부분공간 성질]
$A\in M_{m,n}(K)$에 대하여
\[
\mathcal{N}(A):=\{\mathbf{x}\in K^n: A\mathbf{x}=\mathbf{0}\}
\]
는 $K^n$의 부분공간이다.
\end{theorem}

\begin{proof}
먼저 $A\mathbf{0}=\mathbf{0}$이므로 $\mathbf{0}\in\mathcal{N}(A)$이다.

다음으로 $\mathbf{x},\mathbf{y}\in\mathcal{N}(A)$를 잡으면 $A\mathbf{x}=\mathbf{0}$, $A\mathbf{y}=\mathbf{0}$이므로
\[
A(\mathbf{x}+\mathbf{y})=A\mathbf{x}+A\mathbf{y}=\mathbf{0}+\mathbf{0}=\mathbf{0}.
\]
따라서 $\mathbf{x}+\mathbf{y}\in\mathcal{N}(A)$이다.

또한 임의의 $c\in K$와 $\mathbf{x}\in\mathcal{N}(A)$에 대해
\[
A(c\mathbf{x})=c(A\mathbf{x})=c\mathbf{0}=\mathbf{0}
\]
이므로 $c\mathbf{x}\in\mathcal{N}(A)$이다.

영벡터를 포함하고 덧셈과 스칼라곱에 닫혀 있으므로 $\mathcal{N}(A)$는 $K^n$의 부분공간이다.
\end{proof}

\paragraph{의미 해석.}
동차시스템의 해집합은 단순한 해 목록이 아니라 선형구조를 가진 공간입니다. 이후 장에서 등장하는 차원, 기저, 핵(kernel) 개념이 이 사실 위에서 자연스럽게 연결됩니다.

\begin{example}
다음 시스템을 보자.
\[
\begin{cases}
x_1+2x_2-x_3+3x_4=1,\\
2x_1+x_2+x_3-x_4=0,\\
-x_1+x_2+4x_3+2x_4=3.
\end{cases}
\]
계수행렬과 증강행렬은
\[
A=
\begin{bmatrix}
1&2&-1&3\\
2&1&1&-1\\
-1&1&4&2
\end{bmatrix},
\qquad
(A\mid \mathbf{b})=
\left[
\begin{array}{cccc|c}
1&2&-1&3&1\\
2&1&1&-1&0\\
-1&1&4&2&3
\end{array}
\right]
\]
이다. 같은 $A$에서 $\mathbf{b}=\mathbf{0}$으로 바꾸면 동차시스템이 되며 해집합은 부분공간이 된다.
\end{example}

\subsection*{주의}
\begin{warning}
증강행렬의 마지막 열은 계수열이 아니라 상수열이다. 계수행렬의 열과 같은 방식으로 해석하면 피벗 판정이 틀어질 수 있다.
\end{warning}
\begin{warning}
비동차시스템의 해집합은 일반적으로 부분공간이 아니다. 원점을 지나지 않을 수 있으므로 동차시스템과 같은 공리 검사를 그대로 적용하면 오류가 난다.
\end{warning}

\subsection*{자가진단퀴즈}
\begin{enumerate}[label=\arabic*.]
\item $3$개의 식과 $5$개의 미지수를 갖는 임의의 시스템을 만들고 계수행렬과 증강행렬을 쓰시오.
\item 같은 계수행렬 $A$에 대해 $\mathbf{b}=\mathbf{0}$일 때와 $\mathbf{b}\neq \mathbf{0}$일 때 해집합의 기하적 차이를 설명하시오.
\item $\mathcal{N}(A)$가 공집합이 될 수 없는 이유를 한 줄 증명으로 제시하시오.
\end{enumerate}

\sectiontemplate{기본 행연산과 동치시스템}{Gauss 소거법의 핵심은 계산 편의를 위해 식을 바꾸되 해는 바꾸지 않는다는 점입니다. 이를 보장하는 정확한 규칙이 기본 행연산(elementary row operation)입니다.}{\item 세 가지 기본 행연산을 정확히 정의할 수 있습니다.\item 한 번의 행연산이 해집합을 보존함을 증명할 수 있습니다.\item 유한 번의 행연산으로 얻은 행동치 시스템을 해석할 수 있습니다.}{\item 등식의 동치 변형\item 선형결합의 의미}

\begin{definition}
증강행렬의 행에 대해 다음 연산을 \emph{기본 행연산}이라 한다.
\begin{enumerate}[label=(\alph*)]
\item 두 행을 서로 교환한다.
\item 한 행에 $0$이 아닌 스칼라를 곱한다.
\item 한 행에 다른 행의 스칼라배를 더한다.
\end{enumerate}
\end{definition}

\begin{definition}
행렬 $B$가 행렬 $A$로부터 유한 번의 기본 행연산으로 얻어지면 $A$와 $B$가 \emph{행동치}(row equivalent)라고 한다.
\end{definition}

\paragraph{증명 전략.}
한 번의 행연산에 대해 해집합이 같음을 보일 때는 세 경우를 분리하여 양방향 포함을 확인한다. 특히 (c)형 연산에서는 바뀌지 않은 행을 함께 사용해 원래 식을 복원한다.

\begin{theorem}[한 번의 기본 행연산은 해집합을 보존한다]
연립일차방정식 $S$의 증강행렬에 기본 행연산을 한 번 적용해 얻은 시스템을 $S'$라 하자. 그러면 $S$와 $S'$의 해집합은 같다.
\end{theorem}

\begin{proof}
$S$의 $i$번째 방정식을 $f_i(\mathbf{x})=b_i$라 쓰자.

\textbf{(a) 행 교환}:
두 방정식의 순서만 바뀐다. 방정식의 집합은 동일하므로 해집합이 같다.

\textbf{(b) $i$번째 행에 $c\neq 0$을 곱함}:
새 식은 $c f_i(\mathbf{x})=c b_i$이다.  
$\mathbf{x}$가 원래 식 $f_i(\mathbf{x})=b_i$를 만족하면 새 식도 만족한다.  
반대로 $\mathbf{x}$가 새 식을 만족하면 $c\neq 0$이므로 양변을 $c$로 나누어 원래 식을 얻는다. 다른 행은 변하지 않았으므로 전체 해집합이 같다.

\textbf{(c) $i$번째 행에 $j$번째 행의 $c$배를 더함}:
새 $i$번째 식은
\[
f_i(\mathbf{x})+c f_j(\mathbf{x})=b_i+c b_j
\]
이다. 원래 시스템의 해는 당연히 이 식을 만족한다.  
역으로 새 시스템의 해 $\mathbf{x}$를 잡으면 $j$번째 식 $f_j(\mathbf{x})=b_j$는 그대로 유지되므로
\[
f_i(\mathbf{x})=\bigl(f_i(\mathbf{x})+c f_j(\mathbf{x})\bigr)-c f_j(\mathbf{x})
=(b_i+c b_j)-c b_j=b_i
\]
를 얻는다. 따라서 원래 $i$번째 식도 만족한다.

세 경우 모두에서 해집합이 같으므로 정리가 성립한다.
\end{proof}

\paragraph{의미 해석.}
이 정리는 소거 계산의 정당성 자체입니다. 즉, 우리는 문제를 바꾸는 것이 아니라 같은 해를 갖는 더 쉬운 표현으로 옮기고 있는 것입니다.

\paragraph{증명 전략.}
앞 정리를 반복 적용하는 귀납 구조를 사용한다. 길이 $t$인 행연산 열을 마지막 연산과 앞부분으로 분해하면 증명이 짧아진다.

\begin{theorem}[행동치 시스템의 해집합 일치]
두 시스템의 증강행렬이 행동치이면 두 시스템의 해집합은 같다.
\end{theorem}

\begin{proof}
행동치라는 가정에 의해 한 시스템에서 다른 시스템으로 가는 기본 행연산 열이 존재한다. 그 길이를 $t$라 하자.

$t=1$이면 앞 정리에 의해 해집합이 같다.

$t\ge 2$라고 하자. 처음 $t-1$번 행연산으로 얻은 중간 시스템을 $S_{t-1}$이라 두면, 귀납가정으로 원래 시스템과 $S_{t-1}$의 해집합이 같다. 마지막 1번 행연산으로 $S_t$를 얻는데, 앞 정리에 의해 $S_{t-1}$와 $S_t$의 해집합이 같다. 따라서 원래 시스템과 $S_t$의 해집합도 같다.

귀납법으로 모든 $t$에 대해 성립한다.
\end{proof}

\paragraph{의미 해석.}
실제 계산에서는 여러 번의 행연산을 연속으로 쓰기 때문에, 이 정리가 있어야 알고리즘 전체의 신뢰성이 확보됩니다. 문제 풀이에서 "동치 변형"이라는 말을 안전하게 쓸 수 있는 근거가 됩니다.

\begin{example}
\[
\left[
\begin{array}{ccc|c}
1&-1&2&3\\
2&1&1&4\\
-1&2&-1&-1
\end{array}
\right]
\overset{R_2\leftarrow R_2-2R_1}{\sim}
\left[
\begin{array}{ccc|c}
1&-1&2&3\\
0&3&-3&-2\\
-1&2&-1&-1
\end{array}
\right]
\overset{R_3\leftarrow R_3+R_1}{\sim}
\left[
\begin{array}{ccc|c}
1&-1&2&3\\
0&3&-3&-2\\
0&1&1&2
\end{array}
\right].
\]
세 증강행렬은 서로 행동치이며 같은 해집합을 갖는다.
\end{example}

\subsection*{주의}
\begin{warning}
열연산은 일반적으로 연립일차방정식의 해집합을 보존하지 않는다. 이 장의 해법에서는 행연산만 허용된다.
\end{warning}
\begin{warning}
$R_i\leftarrow 0\cdot R_i$는 기본 행연산이 아니다. 이 연산은 식 정보를 소거해 동치성을 파괴한다.
\end{warning}

\subsection*{자가진단퀴즈}
\begin{enumerate}[label=\arabic*.]
\item 각 기본 행연산의 역연산을 명시하고, 왜 같은 종류의 행연산인지 설명하시오.
\item 행 교환만 여러 번 수행한 경우 해집합이 변하지 않는 이유를 집합 관점에서 서술하시오.
\item $R_i\leftarrow R_i+cR_j$에서 $j$번째 식이 왜 반드시 필요했는지 증명 관점으로 설명하시오.
\end{enumerate}

\sectiontemplate{Gauss--Jordan 소거와 RREF}{행연산의 정당성이 확보되었으므로 이제 계산 표준형을 정합니다. 이 절에서는 사다리꼴과 기약행 사다리꼴을 정의하고, RREF의 존재와 유일성을 완전하게 증명하겠습니다.}{\item REF와 RREF 조건을 정확히 구분할 수 있습니다.\item 모든 행렬이 RREF와 행동치임을 설명할 수 있습니다.\item RREF가 유일하다는 사실을 증명 흐름으로 재현할 수 있습니다.}{\item 기본 행연산\item 랭크의 기초 의미(피벗 개수)}

\begin{definition}
행렬 $U\in M_{m,n}(K)$가 다음을 만족하면 \emph{행 사다리꼴}(row echelon form, REF)이라 한다.
\begin{enumerate}[label=(\roman*)]
\item 영이 아닌 행은 영행보다 위에 있다.
\item 각 영이 아닌 행의 첫 비영원소(leading entry)는 바로 위 행의 leading entry보다 오른쪽에 있다.
\item 각 leading entry 아래의 원소는 모두 $0$이다.
\end{enumerate}
\end{definition}

\begin{definition}
REF 행렬 $R$이 추가로 다음을 만족하면 \emph{기약행 사다리꼴}(reduced row echelon form, RREF)이라 한다.
\begin{enumerate}[label=(\roman*)]
\item 각 leading entry는 $1$이다.
\item leading entry가 있는 열에서는 그 원소를 제외한 모든 원소가 $0$이다.
\end{enumerate}
leading entry가 놓인 열을 \emph{피벗열}(pivot column), 피벗이 없는 열의 변수들을 \emph{자유변수}(free variable)라 한다.
\end{definition}

\paragraph{증명 전략.}
알고리즘을 명시적으로 제시하고 종료성을 보인다. 전진 소거로 REF를 만든 뒤, 후진 소거로 피벗 위 원소를 없애면 RREF가 완성된다.

\begin{theorem}[RREF의 존재]
임의의 $A\in M_{m,n}(K)$에 대해 $A$와 행동치인 RREF 행렬이 존재한다.
\end{theorem}

\begin{proof}
$A$에 대해 다음 절차를 수행한다.

\textbf{1단계(전진 소거)}:
행 인덱스 $i=1$, 열 인덱스 $j=1$로 시작한다.
\begin{enumerate}[label=(\alph*)]
\item $i>m$ 또는 $j>n$이면 종료한다.
\item 열 $j$의 $i$행부터 $m$행까지가 모두 $0$이면 $j\leftarrow j+1$로 바꾸고 반복한다.
\item 그렇지 않으면 어떤 $k\ge i$에 대해 $a_{kj}\neq 0$이 있으므로 $R_i\leftrightarrow R_k$로 교환한다.
\item $R_i\leftarrow a_{ij}^{-1}R_i$로 바꾸어 피벗을 $1$로 만든다.
\item 모든 $\ell>i$에 대해 $R_\ell\leftarrow R_\ell-a_{\ell j}R_i$를 적용해 피벗 아래를 $0$으로 만든다.
\item $i\leftarrow i+1$, $j\leftarrow j+1$로 갱신하고 반복한다.
\end{enumerate}

이 과정에서 각 반복마다 $i$ 또는 $j$가 증가하므로 반복 횟수는 유한하고, 종료 시점 행렬 $U$는 REF 조건을 만족한다.

\textbf{2단계(후진 소거)}:
$U$의 마지막 피벗행부터 첫 피벗행까지 거꾸로 올라오며, 각 피벗열에서 위쪽 원소를 행연산
\[
R_\ell\leftarrow R_\ell-u_{\ell p}R_{i}
\]
로 소거한다($p$는 해당 피벗열, $\ell<i$). 유한 개의 행과 열에서 수행되므로 이 단계도 유한 번에 끝난다. 종료 행렬을 $R$이라 하면 REF 성질은 유지되고 각 피벗열에서 피벗 위아래가 모두 $0$, 피벗은 $1$이므로 $R$은 RREF이다.

모든 단계가 기본 행연산으로 이루어졌으므로 $R$은 $A$와 행동치이다.
\end{proof}

\paragraph{의미 해석.}
이 정리는 소거법이 단순한 요령이 아니라 항상 종료하는 알고리즘임을 보장합니다. 따라서 손계산, 코드 구현, 이론 증명에서 동일한 표준형을 사용할 수 있습니다.

\paragraph{증명 전략.}
유일성은 두 가지 불변량을 사용해 보인다. 먼저 ``앞에서 $k$개 열의 랭크''를 비교해 피벗 위치가 같음을 보이고, 이어서 비피벗열의 열결합 계수가 같음을 보여 전체 행렬 일치를 결론낸다.

\begin{theorem}[RREF의 유일성]
행렬 $A$와 행동치인 RREF는 하나뿐이다.
\end{theorem}

\begin{proof}
$A$와 행동치인 두 RREF를 $R,S\in M_{m,n}(K)$라 하자. $R=EA$, $S=FA$인 가역행렬 $E,F\in M_m(K)$가 존재한다.

행렬 $M$의 앞에서 $k$개 열로 이루어진 부분행렬을 $M^{(k)}$라 쓰자. 그러면
\[
R^{(k)}=E A^{(k)},\qquad S^{(k)}=F A^{(k)}
\]
이므로 가역행렬의 왼쪽 곱에 의한 랭크 불변성에 의해
\[
\operatorname{rank}(R^{(k)})=\operatorname{rank}(A^{(k)})=\operatorname{rank}(S^{(k)})
\quad (k=1,\dots,n).
\]
RREF에서 $\operatorname{rank}(M^{(k)})$는 앞 $k$개 열 안에 있는 피벗열 개수와 같다. 따라서 모든 $k$에서 $R$과 $S$의 ``앞 $k$개 열 피벗 개수''가 같고, 결과적으로 피벗열 위치가 열별로 일치한다. 공통 피벗열을
\[
p_1<\cdots<p_r
\]
라 하자.

이제 행렬 $M$의 $j$번째 열벡터를 $C_j(M)$로 표기한다. $j\notin\{p_1,\dots,p_r\}$를 비피벗열이라 하자. RREF 성질상 $R$의 피벗열은 표준기저열이므로
\[
C_j(R)=\sum_{i=1}^r r_{ij}\,C_{p_i}(R)
\]
가 성립한다($r_{ij}$는 $R$의 $(i,j)$성분). $R$과 $S$는 서로 행동치이므로 $S=GR$인 가역행렬 $G$가 존재한다. 위 식에 $G$를 곱하면
\[
C_j(S)=\sum_{i=1}^r r_{ij}\,C_{p_i}(S).
\]
그런데 $S$도 RREF이므로 $C_{p_i}(S)$ 역시 $i$번째 피벗행에만 $1$이 있고 나머지가 $0$인 표준기저열이다. 따라서 $j$번째 열의 피벗행 성분은
\[
s_{ij}=r_{ij}\quad (i=1,\dots,r)
\]
가 된다. 즉 모든 비피벗열에서 피벗행 성분이 일치한다.

피벗열은 이미 동일한 표준기저열이고, 랭크가 $r$이므로 두 행렬의 나머지 $m-r$개 행은 모두 영행이다. 따라서 모든 성분이 일치하여 $R=S$이다.
\end{proof}

\paragraph{의미 해석.}
소거 경로가 달라도 최종 RREF가 같다는 사실은 계산 검증의 기준을 제공합니다. 서로 다른 풀이가 같은 RREF에 도달하면 정답 신뢰도를 높일 수 있습니다.

\begin{example}
\[
\left[
\begin{array}{cccc|c}
1&2&-1&3&4\\
2&4&-2&6&8\\
1&1&1&1&2
\end{array}
\right]
\sim
\left[
\begin{array}{cccc|c}
1&0&3&-1&0\\
0&1&-2&2&2\\
0&0&0&0&0
\end{array}
\right].
\]
피벗열은 $1,2$열이고 $3,4$열은 자유열이다. 어떤 행연산 순서를 택해도 최종 RREF는 위와 동일하다.
\end{example}

\subsection*{주의}
\begin{warning}
REF와 RREF를 혼동하면 자유변수 판정이 틀어진다. 해를 즉시 읽으려면 RREF 조건(피벗열의 위아래 소거)을 반드시 확인해야 한다.
\end{warning}
\begin{warning}
계산 중 분수 출현을 무리하게 피하면 중간 수가 커져 실수가 늘어날 수 있다. 피벗 정규화 시점을 지나치게 늦추지 않는 것이 안전하다.
\end{warning}

\subsection*{자가진단퀴즈}
\begin{enumerate}[label=\arabic*.]
\item 주어진 REF가 RREF인지 판정하기 위한 체크리스트를 작성하시오.
\item 앞에서 $k$개 열 랭크가 피벗 위치를 결정한다는 사실을 자신의 말로 정리하시오.
\item 임의의 $3\times 4$ 행렬을 하나 정해 서로 다른 두 소거 경로로 RREF를 계산하고 결과 일치를 확인하시오.
\end{enumerate}

\sectiontemplate{해의 존재성, 자유변수, 일반해}{이제 RREF를 해석 도구로 사용합니다. 해가 존재하는지, 몇 개의 자유도가 남는지, 전체 해를 어떤 꼴로 써야 하는지를 한 번에 정리하겠습니다.}{\item 랭크 조건으로 해의 존재성을 판정할 수 있습니다.\item 자유변수 개수를 $n-\operatorname{rank}(A)$로 계산할 수 있습니다.\item 비동차 해집합을 특수해와 동차해의 합으로 표현할 수 있습니다.}{\item RREF 해석\item 동차해집합의 부분공간 성질}

\begin{definition}
$A\in M_{m,n}(K)$의 \emph{랭크}(rank) $\operatorname{rank}(A)$를 $A$의 피벗열 개수로 정의한다. 또한 동차시스템 $A\mathbf{x}=\mathbf{0}$의 해집합을 $\ker A$로 쓴다.
\end{definition}

\paragraph{증명 전략.}
증강행렬을 RREF로 보내 모순행의 존재 여부를 해 존재성과 연결한다. 이어서 피벗열/자유열 분해로 자유매개변수 개수를 바로 계산한다.

\begin{theorem}[Rouch\'e--Capelli와 자유도]
시스템 $A\mathbf{x}=\mathbf{b}$가 해를 가질 필요충분조건은
\[
\operatorname{rank}(A)=\operatorname{rank}(A\mid \mathbf{b})
\]
이다. 해가 존재하면 자유변수 개수는 $n-\operatorname{rank}(A)$이며, 유일해 조건은 $\operatorname{rank}(A)=n$이다.
\end{theorem}

\begin{proof}
증강행렬 $(A\mid\mathbf{b})$를 RREF로 변환하여 $(R\mid\mathbf{c})$를 얻는다. 행동치 정리에 의해 두 시스템의 해집합은 같다.

RREF에서 해가 없다는 것은 어떤 행이
\[
[\,0\ \cdots\ 0\mid d\,],\quad d\neq 0
\]
형태가 되는 것과 동치이다. 이 행은 계수부분에는 영행이지만 증강열 때문에 전체 증강행렬에서는 비영행이므로
\[
\operatorname{rank}(R\mid\mathbf{c})>\operatorname{rank}(R)
\]
를 강제한다. 랭크는 행연산으로 불변이므로
\[
\operatorname{rank}(A\mid\mathbf{b})>\operatorname{rank}(A)
\]
가 된다.

반대로 \(\operatorname{rank}(A\mid\mathbf{b})>\operatorname{rank}(A)\)이면 \((R\mid\mathbf{c})\)에서 마지막 열이 피벗열이 되는 행이 존재하고, 이는 위 모순행과 동일한 의미이므로 시스템은 불능이다. 따라서 해 존재성과 랭크 일치가 서로 동치이다.

이제 해가 존재한다고 가정하자. 계수부분 \(R\)의 피벗열 수를 \(r=\operatorname{rank}(A)\)라 하면, 나머지 \(n-r\)개 열은 자유열이다. 자유열 변수들을 매개변수로 두면 각 피벗변수는 RREF 식에 의해 유일하게 결정된다. 그러므로 자유변수 개수는 정확히 \(n-r\)이다. 유일해는 자유변수가 없는 경우와 동치이므로 \(n-r=0\), 즉 \(r=n\)과 동치이다.
\end{proof}

\paragraph{의미 해석.}
랭크 하나로 ``해의 존재''와 ``해의 개수 구조''를 동시에 읽을 수 있습니다. 계산 결과를 단순 숫자가 아니라 구조 정보로 바꿔 주는 핵심 정리입니다.

\paragraph{증명 전략.}
한 특수해를 기준점으로 잡고 임의의 해와의 차를 보면 동차해가 된다. 반대로 동차해를 특수해에 더하면 다시 비동차해가 된다는 양방향 포함을 확인한다.

\begin{theorem}[비동차 해집합의 평행이동 표현]
$A\mathbf{x}=\mathbf{b}$가 해를 갖고 $\mathbf{x}_p$를 그 한 해라 하자. 그러면 전체 해집합은
\[
\mathbf{x}_p+\ker A=\{\mathbf{x}_p+\mathbf{z}:\mathbf{z}\in\ker A\}
\]
와 같다.
\end{theorem}

\begin{proof}
먼저 $\mathbf{x}$가 $A\mathbf{x}=\mathbf{b}$의 임의의 해라 하자. 그러면
\[
A(\mathbf{x}-\mathbf{x}_p)=A\mathbf{x}-A\mathbf{x}_p=\mathbf{b}-\mathbf{b}=\mathbf{0}
\]
이므로 $\mathbf{x}-\mathbf{x}_p\in\ker A$이다. 따라서 $\mathbf{x}=\mathbf{x}_p+\mathbf{z}$ ($\mathbf{z}\in\ker A$) 꼴로 쓸 수 있다. 즉 해집합은 $\mathbf{x}_p+\ker A$에 포함된다.

반대로 $\mathbf{z}\in\ker A$를 택하면
\[
A(\mathbf{x}_p+\mathbf{z})=A\mathbf{x}_p+A\mathbf{z}=\mathbf{b}+\mathbf{0}=\mathbf{b}
\]
이므로 $\mathbf{x}_p+\mathbf{z}$는 원래 비동차시스템의 해이다. 따라서 $\mathbf{x}_p+\ker A$가 해집합에 포함된다.

양방향 포함으로 두 집합이 같다.
\end{proof}

\paragraph{의미 해석.}
비동차 문제는 ``기준점 하나(특수해) + 방향공간 하나(동차해공간)''로 분해됩니다. 이 관점은 선형대수 전반에서 해석을 단순화하는 표준 틀입니다.

\begin{example}
\[
\begin{cases}
x_1+x_2+x_3=2,\\
2x_1+2x_2+2x_3=4,\\
x_1-x_2+x_3=0
\end{cases}
\]
의 증강행렬을 소거하면
\[
\left[
\begin{array}{ccc|c}
1&0&1&1\\
0&1&0&1\\
0&0&0&0
\end{array}
\right]
\]
을 얻는다. 따라서
\[
x_1=1-t,\quad x_2=1,\quad x_3=t\quad (t\in K).
\]
특수해를 $\mathbf{x}_p=(1,1,0)^T$로 잡으면
\[
\ker A=\operatorname{span}\{(-1,0,1)^T\},
\quad
\mathbf{x}=\mathbf{x}_p+t(-1,0,1)^T
\]
가 되어 위 정리를 그대로 확인할 수 있다.
\end{example}

\subsection*{주의}
\begin{warning}
$\operatorname{rank}(A)=m$이라는 정보만으로는 유일해를 결론낼 수 없다. 유일성은 열 수 $n$과의 비교, 즉 $\operatorname{rank}(A)=n$ 조건이 필요하다.
\end{warning}
\begin{warning}
특수해는 하나로 고정되지 않는다. 서로 다른 특수해의 차가 항상 $\ker A$에 속한다는 점을 이용해 표현의 동치를 확인해야 한다.
\end{warning}

\subsection*{자가진단퀴즈}
\begin{enumerate}[label=\arabic*.]
\item $m<n$인 동차시스템이 항상 비자명해를 갖는 이유를 랭크로 설명하시오.
\item $\operatorname{rank}(A)=\operatorname{rank}(A\mid\mathbf{b})=n$일 때 왜 유일해가 되는지 소거식 형태로 설명하시오.
\item $\mathbf{x}_p+\ker A$ 표현에서 $\mathbf{x}_p$를 다른 특수해로 바꿔도 같은 해집합이 되는 이유를 증명하시오.
\end{enumerate}

\chapterapplicationtemplate{세 갈래 전기회로의 가지전류 $i_1,i_2,i_3$가
\[
\begin{cases}
i_1-i_2-i_3=0,\\
2i_1+i_2=5,\\
-i_1+3i_3=1
\end{cases}
\]
를 만족한다고 하자. 증강행렬을 세워 Gauss--Jordan 소거로 전류를 구하고, 계수의 작은 변화가 있을 때 해의 존재성과 자유도가 랭크 조건으로 어떻게 판정되는지 해석하라. 이 과정은 실제 모델에서 ``식 세우기 $\to$ 소거 $\to$ 구조 해석''의 표준 흐름을 보여 준다.}

\chaptersummarytemplate

\section*{종합 연습문제}

\subsection*{연습문제 A (기초 확인)}
\begin{enumerate}[label=A-\arabic*.]
\item 다음 시스템의 증강행렬을 쓰시오.
\[
\begin{cases}
x_1-2x_2+x_3=1,\\
3x_1+x_2-4x_3=0.
\end{cases}
\]
\item 행렬
\[
\left[
\begin{array}{ccc|c}
1&1&0&2\\
2&3&1&5
\end{array}
\right]
\]
에 대해 $R_2\leftarrow R_2-2R_1$을 수행하시오.
\item 다음 중 기본 행연산이 아닌 것을 모두 고르시오.  
(a) $R_1\leftrightarrow R_3$ \quad (b) $R_2\leftarrow 0\cdot R_2$ \quad (c) $R_1\leftarrow R_1+5R_2$
\item 다음 RREF에서 피벗변수와 자유변수를 구분하시오.
\[
\left[
\begin{array}{cccc|c}
1&0&2&-1&3\\
0&1&-4&5&-2\\
0&0&0&0&0
\end{array}
\right].
\]
\item 동차시스템 $A\mathbf{x}=\mathbf{0}$의 해집합이 부분공간임을 공리 3개(영벡터, 덧셈, 스칼라곱)로 직접 확인하시오.
\item 다음 두 시스템이 행동치인지 판정하시오.
\[
\left[
\begin{array}{cc|c}
1&2&1\\
2&4&2
\end{array}
\right],\qquad
\left[
\begin{array}{cc|c}
1&2&1\\
0&0&0
\end{array}
\right].
\]
\item $A\in M_{3,5}(K)$, $\operatorname{rank}(A)=3$일 때 동차시스템 자유변수 개수를 구하시오.
\end{enumerate}

\subsection*{연습문제 B (개념 연결)}
\begin{enumerate}[label=B-\arabic*.]
\item 한 번의 기본 행연산이 해집합을 보존함을 세 경우로 나누어 완전하게 증명하시오.
\item 행동치가 동치관계(반사성, 대칭성, 추이성)를 이룸을 보이시오.
\item $A\mathbf{x}=\mathbf{b}$의 해 존재 조건이 $\operatorname{rank}(A)=\operatorname{rank}(A\mid\mathbf{b})$와 동치임을 모순행 관점에서 설명하시오.
\item 비동차 해집합이 공집합이 아니면 아핀집합이 됨을 $\mathbf{x}_p+\ker A$ 형태로 설명하시오.
\item 매개변수 $\alpha\in K$에 대해
\[
\begin{cases}
x_1+x_2=1,\\
2x_1+2x_2=\alpha
\end{cases}
\]
의 해 존재성을 $\alpha$에 따라 분류하시오.
\end{enumerate}

\subsection*{연습문제 C (증명/종합)}
\begin{enumerate}[label=C-\arabic*.]
\item 다음 사실을 증명하시오: 임의의 행렬 $A$에서 앞 $k$개 열의 랭크 증가 여부는 피벗열 위치를 결정하며, 이를 이용하면 RREF의 피벗 위치가 유일하다.
\item $A\mathbf{x}=\mathbf{b}$가 가해일 때 임의의 두 특수해 $\mathbf{x}_p,\mathbf{x}_q$에 대해 $\mathbf{x}_p-\mathbf{x}_q\in\ker A$임을 보이고, 해집합 표현이 특수해 선택과 무관함을 증명하시오.
\item 다음 명제를 증명하시오: $A\in M_{m,n}(K)$에서 $\operatorname{rank}(A)=n$이면 모든 $\mathbf{b}\in K^m$에 대해 $A\mathbf{x}=\mathbf{b}$가 가해인 것은 일반적으로 거짓이며, 참이 되기 위한 필요충분조건을 $m,n$과 랭크로 제시하시오.
\end{enumerate}